\documentclass[12pt, titlepage]{article}

\usepackage{amsmath, mathtools}

\usepackage[round]{natbib}
\usepackage{amsfonts}
\usepackage{amssymb}
\usepackage{graphicx}
\usepackage{colortbl}
\usepackage{xr}
\usepackage{hyperref}
\usepackage{longtable}
\usepackage{xfrac}
\usepackage{tabularx}
\usepackage{float}
\usepackage{siunitx}
\usepackage{booktabs}
\usepackage{multirow}
\usepackage[section]{placeins}
\usepackage{caption}
\usepackage{fullpage}

\hypersetup{
bookmarks=true,     % show bookmarks bar?
colorlinks=true,       % false: boxed links; true: colored links
linkcolor=red,          % color of internal links (change box color with linkbordercolor)
citecolor=blue,      % color of links to bibliography
filecolor=magenta,  % color of file links
urlcolor=cyan          % color of external links
}

\usepackage{array}

\externaldocument{../../SRS/SRS}

\input{../../Comments.text}
\input{../../Common.text}

\begin{document}

\title{Module Interface Specification for \progname{}}

\author{\authname}

\date{\today}

\maketitle

\pagenumbering{roman}

\section{Revision History}

\begin{tabularx}{\textwidth}{p{3cm}p{2cm}X}
\toprule {\bf Date} & {\bf Version} & {\bf Notes}\\
\midrule
Date 1 & 1.0 & Notes\\
Date 2 & 1.1 & Notes\\
\bottomrule
\end{tabularx}

\section{Symbols, Abbreviations and Acronyms}

See SRS Documentation at \wss{give url}

\wss{Also add any additional symbols, abbreviations or acronyms}

\tableofcontents

\pagenumbering{arabic}

\section{Introduction}

The following document details the Module Interface Specifications for
\wss{Fill in your project name and description}

Complementary documents include the System Requirement Specifications
and Module Guide.  The full documentation and implementation can be
found at \url{...}.  \wss{provide the url for your repo}

\section{Notation}

\wss{You should describe your notation.  You can use what is below as
  a starting point.}

The structure of the MIS for modules comes from \citet{HoffmanAndStrooper1995},
with the addition that template modules have been adapted from
\cite{GhezziEtAl2003}.  The mathematical notation comes from Chapter 3 of
\citet{HoffmanAndStrooper1995}.  For instance, the symbol := is used for a
multiple assignment statement and conditional rules follow the form $(c_1
\Rightarrow r_1 | c_2 \Rightarrow r_2 | ... | c_n \Rightarrow r_n )$.

The following table summarizes the primitive data types used by \progname. 

\begin{center}
\renewcommand{\arraystretch}{1.2}
\noindent 
\begin{tabular}{l l p{7.5cm}} 
\toprule 
\textbf{Data Type} & \textbf{Notation} & \textbf{Description}\\ 
\midrule
character & char & a single symbol or digit\\
integer & $\mathbb{Z}$ & a number without a fractional component in (-$\infty$, $\infty$) \\
natural number & $\mathbb{N}$ & a number without a fractional component in [1, $\infty$) \\
real & $\mathbb{R}$ & any number in (-$\infty$, $\infty$)\\
\bottomrule
\end{tabular} 
\end{center}

\noindent
The specification of \progname \ uses some derived data types: sequences, strings, and
tuples. Sequences are lists filled with elements of the same data type. Strings
are sequences of characters. Tuples contain a list of values, potentially of
different types. In addition, \progname \ uses functions, which
are defined by the data types of their inputs and outputs. Local functions are
described by giving their type signature followed by their specification.

\section{Module Decomposition}

The following table is taken directly from the Module Guide document for this project.

\begin{table}[H]
\centering
\begin{tabular}{p{0.3\textwidth} p{0.6\textwidth}}
\toprule
\textbf{Level 1} & \textbf{Level 2}\\
\midrule

Hardware-Hiding Module & Fitness Integration Module\\
\midrule

\multirow{14}{0.3\textwidth}{Behaviour-Hiding Module} & Persistence Module\\
& User and Onboarding Module\\
& Activity Module\\
& Goals Module\\
& Currency Ledger Module\\
& Asset Module\\
& Transaction Order Module\\
& Portfolio Module\\
& Leaderboard and Friends Module\\
& Progress Reporting Module\\
& UI Screens Module\\
& Notification Module\\
& Market Simulation Module\\
& API Client Module\\
\midrule

Software Decision Module & Authentication and Security Module\\
\bottomrule

\end{tabular}
\caption{Module Hierarchy}
\label{TblMH}
\end{table}

\section{MIS of Fitness Integration Module} \label{misFI}

This module is M1 in the Module Guide.

\subsection{Module}

FitnessIntegration

\subsection{Uses}

\wss{Modules this module uses}

\subsection{Syntax}

\subsubsection{Exported Constants}

\wss{Exported constants, or None}

\subsubsection{Exported Access Programs}

\begin{center}
\begin{tabular}{p{2cm} p{4cm} p{4cm} p{2cm}}
\hline
\textbf{Name} & \textbf{In} & \textbf{Out} & \textbf{Exceptions} \\
\hline
\wss{accessProg} & - & - & - \\
\hline
\end{tabular}
\end{center}

\subsection{Semantics}

\subsubsection{State Variables}

\wss{State variables, or None}

\subsubsection{Environment Variables}

\wss{Environment variables, or None}

\subsubsection{Assumptions}

\wss{Assumptions, or None}

\subsubsection{Access Routine Semantics}

\noindent \wss{accessProg}():
\begin{itemize}
\item transition: \wss{if appropriate}
\item output: \wss{if appropriate}
\item exception: \wss{if appropriate}
\end{itemize}

\subsubsection{Local Functions}

\wss{Local functions, or None}

\section{MIS of Persistence Module} \label{misPer}

This module is M2 in the Module Guide.

\subsection{Module}

Persistence

\subsection{Uses}

\wss{Modules this module uses}

\subsection{Syntax}

\subsubsection{Exported Constants}

\wss{Exported constants, or None}

\subsubsection{Exported Access Programs}

\begin{center}
\begin{tabular}{p{2cm} p{4cm} p{4cm} p{2cm}}
\hline
\textbf{Name} & \textbf{In} & \textbf{Out} & \textbf{Exceptions} \\
\hline
\wss{accessProg} & - & - & - \\
\hline
\end{tabular}
\end{center}

\subsection{Semantics}

\subsubsection{State Variables}

\wss{State variables, or None}

\subsubsection{Environment Variables}

\wss{Environment variables, or None}

\subsubsection{Assumptions}

\wss{Assumptions, or None}

\subsubsection{Access Routine Semantics}

\noindent \wss{accessProg}():
\begin{itemize}
\item transition: \wss{if appropriate}
\item output: \wss{if appropriate}
\item exception: \wss{if appropriate}
\end{itemize}

\subsubsection{Local Functions}

\wss{Local functions, or None}

\section{MIS of User and Onboarding Module} \label{misUser}

This module is M3 in the Module Guide.

\subsection{Module}

User

\subsection{Uses}

\wss{Modules this module uses}

\subsection{Syntax}

\subsubsection{Exported Constants}

\wss{Exported constants, or None}

\subsubsection{Exported Access Programs}

\begin{center}
\begin{tabular}{p{2cm} p{4cm} p{4cm} p{2cm}}
\hline
\textbf{Name} & \textbf{In} & \textbf{Out} & \textbf{Exceptions} \\
\hline
\wss{accessProg} & - & - & - \\
\hline
\end{tabular}
\end{center}

\subsection{Semantics}

\subsubsection{State Variables}

\wss{State variables, or None}

\subsubsection{Environment Variables}

\wss{Environment variables, or None}

\subsubsection{Assumptions}

\wss{Assumptions, or None}

\subsubsection{Access Routine Semantics}

\noindent \wss{accessProg}():
\begin{itemize}
\item transition: \wss{if appropriate}
\item output: \wss{if appropriate}
\item exception: \wss{if appropriate}
\end{itemize}

\subsubsection{Local Functions}

\wss{Local functions, or None}

\section{MIS of Activity Module} \label{misAct}

This module is M4 in the Module Guide.

\subsection{Module}

Activity

\subsection{Uses}

\wss{Modules this module uses}

\subsection{Syntax}

\subsubsection{Exported Constants}

\wss{Exported constants, or None}

\subsubsection{Exported Access Programs}

\begin{center}
\begin{tabular}{p{2cm} p{4cm} p{4cm} p{2cm}}
\hline
\textbf{Name} & \textbf{In} & \textbf{Out} & \textbf{Exceptions} \\
\hline
\wss{accessProg} & - & - & - \\
\hline
\end{tabular}
\end{center}

\subsection{Semantics}

\subsubsection{State Variables}

\wss{State variables, or None}

\subsubsection{Environment Variables}

\wss{Environment variables, or None}

\subsubsection{Assumptions}

\wss{Assumptions, or None}

\subsubsection{Access Routine Semantics}

\noindent \wss{accessProg}():
\begin{itemize}
\item transition: \wss{if appropriate}
\item output: \wss{if appropriate}
\item exception: \wss{if appropriate}
\end{itemize}

\subsubsection{Local Functions}

\wss{Local functions, or None}

\section{MIS of Goals Module} \label{misGoal}

This module is M5 in the Module Guide.

\subsection{Module}

Goals

\subsection{Uses}

\wss{Modules this module uses}

\subsection{Syntax}

\subsubsection{Exported Constants}

\wss{Exported constants, or None}

\subsubsection{Exported Access Programs}

\begin{center}
\begin{tabular}{p{2cm} p{4cm} p{4cm} p{2cm}}
\hline
\textbf{Name} & \textbf{In} & \textbf{Out} & \textbf{Exceptions} \\
\hline
\wss{accessProg} & - & - & - \\
\hline
\end{tabular}
\end{center}

\subsection{Semantics}

\subsubsection{State Variables}

\wss{State variables, or None}

\subsubsection{Environment Variables}

\wss{Environment variables, or None}

\subsubsection{Assumptions}

\wss{Assumptions, or None}

\subsubsection{Access Routine Semantics}

\noindent \wss{accessProg}():
\begin{itemize}
\item transition: \wss{if appropriate}
\item output: \wss{if appropriate}
\item exception: \wss{if appropriate}
\end{itemize}

\subsubsection{Local Functions}

\wss{Local functions, or None}

\section{MIS of Currency Ledger Module} \label{misCur}

This module is M6 in the Module Guide.

\subsection{Module}

CurrencyLedger

\subsection{Uses}

\wss{Modules this module uses}

\subsection{Syntax}

\subsubsection{Exported Constants}

\wss{Exported constants, or None}

\subsubsection{Exported Access Programs}

\begin{center}
\begin{tabular}{p{2cm} p{4cm} p{4cm} p{2cm}}
\hline
\textbf{Name} & \textbf{In} & \textbf{Out} & \textbf{Exceptions} \\
\hline
\wss{accessProg} & - & - & - \\
\hline
\end{tabular}
\end{center}

\subsection{Semantics}

\subsubsection{State Variables}

\wss{State variables, or None}

\subsubsection{Environment Variables}

\wss{Environment variables, or None}

\subsubsection{Assumptions}

\wss{Assumptions, or None}

\subsubsection{Access Routine Semantics}

\noindent \wss{accessProg}():
\begin{itemize}
\item transition: \wss{if appropriate}
\item output: \wss{if appropriate}
\item exception: \wss{if appropriate}
\end{itemize}

\subsubsection{Local Functions}

\wss{Local functions, or None}

\section{MIS of Asset Module} \label{misAsset}

This module is M7 in the Module Guide.

\subsection{Module}

Asset

\subsection{Uses}

\wss{Modules this module uses}

\subsection{Syntax}

\subsubsection{Exported Constants}

\wss{Exported constants, or None}

\subsubsection{Exported Access Programs}

\begin{center}
\begin{tabular}{p{2cm} p{4cm} p{4cm} p{2cm}}
\hline
\textbf{Name} & \textbf{In} & \textbf{Out} & \textbf{Exceptions} \\
\hline
\wss{accessProg} & - & - & - \\
\hline
\end{tabular}
\end{center}

\subsection{Semantics}

\subsubsection{State Variables}

\wss{State variables, or None}

\subsubsection{Environment Variables}

\wss{Environment variables, or None}

\subsubsection{Assumptions}

\wss{Assumptions, or None}

\subsubsection{Access Routine Semantics}

\noindent \wss{accessProg}():
\begin{itemize}
\item transition: \wss{if appropriate}
\item output: \wss{if appropriate}
\item exception: \wss{if appropriate}
\end{itemize}

\subsubsection{Local Functions}

\wss{Local functions, or None}

\section{MIS of Transaction Order Module} \label{misPurch}

This module is M8 in the Module Guide.

\subsection{Module}

TransactionOrder

\subsection{Uses}

\wss{Modules this module uses}

\subsection{Syntax}

\subsubsection{Exported Constants}

\wss{Exported constants, or None}

\subsubsection{Exported Access Programs}

\begin{center}
\begin{tabular}{p{2cm} p{4cm} p{4cm} p{2cm}}
\hline
\textbf{Name} & \textbf{In} & \textbf{Out} & \textbf{Exceptions} \\
\hline
\wss{accessProg} & - & - & - \\
\hline
\end{tabular}
\end{center}

\subsection{Semantics}

\subsubsection{State Variables}

\wss{State variables, or None}

\subsubsection{Environment Variables}

\wss{Environment variables, or None}

\subsubsection{Assumptions}

\wss{Assumptions, or None}

\subsubsection{Access Routine Semantics}

\noindent \wss{accessProg}():
\begin{itemize}
\item transition: \wss{if appropriate}
\item output: \wss{if appropriate}
\item exception: \wss{if appropriate}
\end{itemize}

\subsubsection{Local Functions}

\wss{Local functions, or None}

\section{MIS of Portfolio Module} \label{misPort}

This module is M9 in the Module Guide.

\subsection{Module}

Portfolio

\subsection{Uses}

\wss{Modules this module uses}

\subsection{Syntax}

\subsubsection{Exported Constants}

\wss{Exported constants, or None}

\subsubsection{Exported Access Programs}

\begin{center}
\begin{tabular}{p{2cm} p{4cm} p{4cm} p{2cm}}
\hline
\textbf{Name} & \textbf{In} & \textbf{Out} & \textbf{Exceptions} \\
\hline
\wss{accessProg} & - & - & - \\
\hline
\end{tabular}
\end{center}

\subsection{Semantics}

\subsubsection{State Variables}

\wss{State variables, or None}

\subsubsection{Environment Variables}

\wss{Environment variables, or None}

\subsubsection{Assumptions}

\wss{Assumptions, or None}

\subsubsection{Access Routine Semantics}

\noindent \wss{accessProg}():
\begin{itemize}
\item transition: \wss{if appropriate}
\item output: \wss{if appropriate}
\item exception: \wss{if appropriate}
\end{itemize}

\subsubsection{Local Functions}

\wss{Local functions, or None}

\section{MIS of Leaderboard and Friends Module} \label{misLead}

This module is M10 in the Module Guide.

\subsection{Module}

Leaderboard

\subsection{Uses}

\wss{Modules this module uses}

\subsection{Syntax}

\subsubsection{Exported Constants}

\wss{Exported constants, or None}

\subsubsection{Exported Access Programs}

\begin{center}
\begin{tabular}{p{2cm} p{4cm} p{4cm} p{2cm}}
\hline
\textbf{Name} & \textbf{In} & \textbf{Out} & \textbf{Exceptions} \\
\hline
\wss{accessProg} & - & - & - \\
\hline
\end{tabular}
\end{center}

\subsection{Semantics}

\subsubsection{State Variables}

\wss{State variables, or None}

\subsubsection{Environment Variables}

\wss{Environment variables, or None}

\subsubsection{Assumptions}

\wss{Assumptions, or None}

\subsubsection{Access Routine Semantics}

\noindent \wss{accessProg}():
\begin{itemize}
\item transition: \wss{if appropriate}
\item output: \wss{if appropriate}
\item exception: \wss{if appropriate}
\end{itemize}

\subsubsection{Local Functions}

\wss{Local functions, or None}

\section{MIS of Progress Reporting Module} \label{misProg}

This module is M11 in the Module Guide.

\subsection{Module}

ProgressReport

\subsection{Uses}

\wss{Modules this module uses}

\subsection{Syntax}

\subsubsection{Exported Constants}

\wss{Exported constants, or None}

\subsubsection{Exported Access Programs}

\begin{center}
\begin{tabular}{p{2cm} p{4cm} p{4cm} p{2cm}}
\hline
\textbf{Name} & \textbf{In} & \textbf{Out} & \textbf{Exceptions} \\
\hline
\wss{accessProg} & - & - & - \\
\hline
\end{tabular}
\end{center}

\subsection{Semantics}

\subsubsection{State Variables}

\wss{State variables, or None}

\subsubsection{Environment Variables}

\wss{Environment variables, or None}

\subsubsection{Assumptions}

\wss{Assumptions, or None}

\subsubsection{Access Routine Semantics}

\noindent \wss{accessProg}():
\begin{itemize}
\item transition: \wss{if appropriate}
\item output: \wss{if appropriate}
\item exception: \wss{if appropriate}
\end{itemize}

\subsubsection{Local Functions}

\wss{Local functions, or None}

\section{MIS of UI Screens Module} \label{misUI}

This module is M12 in the Module Guide.

\subsection{Module}

UIScreens

\subsection{Uses}

\wss{Modules this module uses}

\subsection{Syntax}

\subsubsection{Exported Constants}

\wss{Exported constants, or None}

\subsubsection{Exported Access Programs}

\begin{center}
\begin{tabular}{p{2cm} p{4cm} p{4cm} p{2cm}}
\hline
\textbf{Name} & \textbf{In} & \textbf{Out} & \textbf{Exceptions} \\
\hline
\wss{accessProg} & - & - & - \\
\hline
\end{tabular}
\end{center}

\subsection{Semantics}

\subsubsection{State Variables}

\wss{State variables, or None}

\subsubsection{Environment Variables}

\wss{Environment variables, or None}

\subsubsection{Assumptions}

\wss{Assumptions, or None}

\subsubsection{Access Routine Semantics}

\noindent \wss{accessProg}():
\begin{itemize}
\item transition: \wss{if appropriate}
\item output: \wss{if appropriate}
\item exception: \wss{if appropriate}
\end{itemize}

\subsubsection{Local Functions}

\wss{Local functions, or None}

\section{MIS of Notification Module} \label{misNoti}

This module is M13 in the Module Guide.

\subsection{Module}

Notification

\subsection{Uses}

\wss{Modules this module uses}

\subsection{Syntax}

\subsubsection{Exported Constants}

\wss{Exported constants, or None}

\subsubsection{Exported Access Programs}

\begin{center}
\begin{tabular}{p{2cm} p{4cm} p{4cm} p{2cm}}
\hline
\textbf{Name} & \textbf{In} & \textbf{Out} & \textbf{Exceptions} \\
\hline
\wss{accessProg} & - & - & - \\
\hline
\end{tabular}
\end{center}

\subsection{Semantics}

\subsubsection{State Variables}

\wss{State variables, or None}

\subsubsection{Environment Variables}

\wss{Environment variables, or None}

\subsubsection{Assumptions}

\wss{Assumptions, or None}

\subsubsection{Access Routine Semantics}

\noindent \wss{accessProg}():
\begin{itemize}
\item transition: \wss{if appropriate}
\item output: \wss{if appropriate}
\item exception: \wss{if appropriate}
\end{itemize}

\subsubsection{Local Functions}

\wss{Local functions, or None}

\section{MIS of Market Simulation Module} \label{misMkt}

This module is M14 in the Module Guide.

\subsection{Module}

MarketSimulation

\subsection{Uses}

\wss{Modules this module uses}

\subsection{Syntax}

\subsubsection{Exported Constants}

\wss{Exported constants, or None}

\subsubsection{Exported Access Programs}

\begin{center}
\begin{tabular}{p{2cm} p{4cm} p{4cm} p{2cm}}
\hline
\textbf{Name} & \textbf{In} & \textbf{Out} & \textbf{Exceptions} \\
\hline
\wss{accessProg} & - & - & - \\
\hline
\end{tabular}
\end{center}

\subsection{Semantics}

\subsubsection{State Variables}

\wss{State variables, or None}

\subsubsection{Environment Variables}

\wss{Environment variables, or None}

\subsubsection{Assumptions}

\wss{Assumptions, or None}

\subsubsection{Access Routine Semantics}

\noindent \wss{accessProg}():
\begin{itemize}
\item transition: \wss{if appropriate}
\item output: \wss{if appropriate}
\item exception: \wss{if appropriate}
\end{itemize}

\subsubsection{Local Functions}

\wss{Local functions, or None}

\section{MIS of API Client Module} \label{misAPI}

This module is M15 in the Module Guide.

\subsection{Module}

ApiClient

\subsection{Uses}

\wss{Modules this module uses}

\subsection{Syntax}

\subsubsection{Exported Constants}

\wss{Exported constants, or None}

\subsubsection{Exported Access Programs}

\begin{center}
\begin{tabular}{p{2cm} p{4cm} p{4cm} p{2cm}}
\hline
\textbf{Name} & \textbf{In} & \textbf{Out} & \textbf{Exceptions} \\
\hline
\wss{accessProg} & - & - & - \\
\hline
\end{tabular}
\end{center}

\subsection{Semantics}

\subsubsection{State Variables}

\wss{State variables, or None}

\subsubsection{Environment Variables}

\wss{Environment variables, or None}

\subsubsection{Assumptions}

\wss{Assumptions, or None}

\subsubsection{Access Routine Semantics}

\noindent \wss{accessProg}():
\begin{itemize}
\item transition: \wss{if appropriate}
\item output: \wss{if appropriate}
\item exception: \wss{if appropriate}
\end{itemize}

\subsubsection{Local Functions}

\wss{Local functions, or None}

\section{MIS of Authentication and Security Module} \label{misAuth}

This module is M16 in the Module Guide.

\subsection{Module}

Auth

\subsection{Uses}

\begin{itemize}
\item Persistence Module (Section~\ref{misPer}): stores user records, sessions,
  sign-in lockout state, multi-factor authentication challenges and password
  reset tokens.
\end{itemize}

\noindent Auth runs on the backend only. Storing tokens on the phone and
attaching them to requests is the responsibility of the API Client Module
(Section~\ref{misAPI}).

\subsection{Syntax}

\subsubsection{Exported Constants}

The following constants are defined in the SRS.

\begin{center}
\renewcommand{\arraystretch}{1.2}
\begin{tabular}{l l l}
\toprule
\textbf{Constant} & \textbf{Value} & \textbf{SRS Requirement}\\
\midrule
\textsf{AccessExpiry} & 60 minutes & SEC-ACC-01\\
\textsf{SessionTimeout} & 30 days without use & SEC-ACC-05\\
\textsf{MinPasswordLength} & 8 characters & SEC-ACC-04\\
\textsf{LockoutAttempts} & 5 consecutive failures & SEC-ACC-04\\
\textsf{LockoutDuration} & 15 minutes & SEC-ACC-04\\
\textsf{RateLimit} & 60 requests per minute & SEC-IMM-02\\
\bottomrule
\end{tabular}
\end{center}

\noindent Two further constants are local to this module and are not exported:
\textsf{MfaCodeExpiry} = 10 minutes, the time a multi-factor authentication
code remains valid, and \textsf{ResetTokenExpiry} = 30 minutes, the time a
password reset link remains valid.

\subsubsection{Exported Types}

\begin{itemize}
\item UserID = $\mathbb{N}$
\item SessionID = $\mathbb{N}$
\item DeviceID = String, an identifier the client generates when the app is
  installed
\item Token = String
\item Time = $\mathbb{R}$, seconds since the Unix epoch
\item Role = $\{\mathrm{STANDARD}, \mathrm{ADMIN}\}$
\item Principal = tuple of ($\mathit{user}$: UserID, $\mathit{role}$: Role,
  $\mathit{session}$: SessionID, $\mathit{device}$: DeviceID)
\item TokenPair = tuple of ($\mathit{access}$: Token, $\mathit{refresh}$:
  Token, $\mathit{accessExpires}$: Time)
\item MfaChallenge = tuple of ($\mathit{challengeId}$: String)
\item SignInResult = TokenPair $\cup$ MfaChallenge
\end{itemize}

\subsubsection{Exported Access Programs}

\begin{center}
\renewcommand{\arraystretch}{1.2}
\begin{tabular}{>{\raggedright\arraybackslash}p{3.3cm}
  >{\raggedright\arraybackslash}p{4.6cm}
  >{\raggedright\arraybackslash}p{2.4cm}
  >{\raggedright\arraybackslash}p{4cm}}
\toprule
\textbf{Name} & \textbf{In} & \textbf{Out} & \textbf{Exceptions} \\
\midrule
hashPassword & $\mathit{pw}$: String & String & WeakPassword \\
signIn & $\mathit{email}$: String, $\mathit{pw}$: String, $\mathit{dev}$:
  DeviceID, $\mathit{ip}$: String & SignInResult & RateLimited,
  UnknownAccount, AccountLocked, InvalidCredentials \\
verifyMfa & $\mathit{cid}$: String, $\mathit{code}$: String, $\mathit{ip}$:
  String & TokenPair & RateLimited, InvalidCode \\
refresh & $\mathit{rt}$: Token & TokenPair & Unauthenticated \\
authenticate & $\mathit{tok}$: Token, $\mathit{ip}$: String & Principal &
  Unauthenticated, RateLimited \\
authorize & $p$: Principal, $\mathit{owner}$: UserID & - & Forbidden \\
requireRole & $p$: Principal, $r$: Role & - & Forbidden \\
confirmPassword & $p$: Principal, $\mathit{pw}$: String & - &
  InvalidCredentials \\
signOut & $p$: Principal & - & - \\
endAllSessions & $u$: UserID & - & - \\
setMfaEnabled & $p$: Principal, $\mathit{on}$: $\mathbb{B}$ & - & - \\
requestPassword\-Reset & $\mathit{email}$: String, $\mathit{ip}$: String & - &
  RateLimited \\
resetPassword & $t$: Token, $\mathit{pw}$: String & - & InvalidResetToken,
  WeakPassword \\
\bottomrule
\end{tabular}
\end{center}

\noindent The backend reports each exception to the client with the following
HTTP status: Unauthenticated, UnknownAccount, InvalidCredentials and
InvalidCode with 401; Forbidden with 403; AccountLocked with 423; WeakPassword
and InvalidResetToken with 400; and RateLimited with 429.

\subsection{Semantics}

\subsubsection{State Variables}

\begin{itemize}
\item $\mathit{sessions}$: set of Session, where Session = tuple of
  ($\mathit{id}$: SessionID, $\mathit{user}$: UserID, $\mathit{device}$:
  DeviceID, $\mathit{refreshHash}$: String, $\mathit{lastUsed}$: Time,
  $\mathit{revoked}$: $\mathbb{B}$)
\item $\mathit{failedAttempts}$: UserID $\rightarrow \mathbb{Z}$, the number
  of consecutive failed sign-in attempts (initially 0)
\item $\mathit{lockedUntil}$: UserID $\rightarrow$ Time, the time until which
  sign-in is locked (initially 0)
\item $\mathit{mfaEnabled}$: UserID $\rightarrow \mathbb{B}$ (initially false)
\item $\mathit{challenges}$: set of tuples of ($\mathit{id}$: String,
  $\mathit{user}$: UserID, $\mathit{device}$: DeviceID, $\mathit{codeHash}$:
  String, $\mathit{expires}$: Time)
\item $\mathit{resetTokens}$: set of tuples of ($\mathit{tokenHash}$: String,
  $\mathit{user}$: UserID, $\mathit{expires}$: Time)
\item $\mathit{requestLog}$: String $\rightarrow$ sequence of Time, the times
  of recent requests, keyed by client IP address or by user ID
\end{itemize}

\noindent All state variables except $\mathit{requestLog}$ are stored through
the Persistence Module. $\mathit{requestLog}$ is held in memory.

\subsubsection{Environment Variables}

\begin{itemize}
\item $\mathit{clock}$: the system clock; $\mathit{now}()$ returns the
  current Time.
\item $\mathit{signingKey}$: the secret key used to sign access tokens. It is
  read from the deployment environment and is never stored in the project
  repository (SEC-ACC-03).
\item $\mathit{mailServer}$: the outgoing email server used to send
  multi-factor authentication codes and password reset links.
\item $\mathit{securityLog}$: the structured log to which security events are
  written (SEC-AUD-02).
\item $\mathit{CommonPasswords}$: set of String, a list of common passwords
  read from a file when the backend starts (SEC-ACC-04).
\end{itemize}

\subsubsection{Assumptions}

\begin{itemize}
\item Client--server communication uses TLS 1.3, configured when the backend
  is deployed (SEC-PRV-01). This module does not enforce it.
\item The User and Onboarding Module ensures that email addresses are unique.
\item Each client sends the same DeviceID on every request from a given
  installation.
\item A single backend instance is running, so the in-memory
  $\mathit{requestLog}$ counts every request. Running several instances would
  require a shared store.
\item authenticate is called before any other backend access program on every
  protected request. Registration, signIn, verifyMfa, refresh,
  requestPasswordReset, resetPassword and the health check are the only
  unprotected requests.
\end{itemize}

\subsubsection{Access Routine Semantics}

\noindent Each call to signIn, verifyMfa or requestPasswordReset first appends
$\mathit{now}()$ to $\mathit{requestLog}(\mathit{ip})$. When a token passes
verification in authenticate, $\mathit{now}()$ is first appended to
$\mathit{requestLog}(p.\mathit{user})$. Every exception raised by this module is
written to $\mathit{securityLog}$ with $\mathit{log}$, together with every
account lockout. In the routines below, an exception leaves the state
unchanged unless a transition for that exception is given.

\medskip

\noindent hashPassword($\mathit{pw}$):
\begin{itemize}
\item output: $\mathit{out} := \mathit{hash}(\mathit{pw})$
\item exception: $\mathit{isWeak}(\mathit{pw}) \Rightarrow
  \mathrm{WeakPassword}$
\end{itemize}

\noindent signIn($\mathit{email}, \mathit{pw}, \mathit{dev}, \mathit{ip}$):
\begin{itemize}
\item Let $a = \mathit{findUser}(\mathit{email})$, $\mathit{ok} =
  \mathit{matches}(\mathit{pw}, a.\mathit{passwordHash})$, $\mathit{rt} =
  \mathit{newSecret}()$, $c = \mathit{newCode}()$, and $s =
  (\mathit{newId}(), a.\mathit{id}, \mathit{dev}, \mathit{digest}(\mathit{rt}),
  \mathit{now}(), \mathrm{false})$.
\item transition:
  \begin{itemize}
  \item $\mathit{ok} \wedge \mathit{mfaEnabled}(a.\mathit{id}) \Rightarrow
    \mathit{failedAttempts}(a.\mathit{id}) := 0$. A new challenge with id
    $\mathit{newId}()$, user $a.\mathit{id}$, device $\mathit{dev}$, code hash
    $\mathit{hash}(c)$ and expiry $\mathit{now}() + \textsf{MfaCodeExpiry}$ is
    added to $\mathit{challenges}$, and $c$ is emailed to $a.\mathit{email}$
    through $\mathit{mailServer}$.
  \item $\mathit{ok} \wedge \neg \mathit{mfaEnabled}(a.\mathit{id}) \Rightarrow
    \mathit{failedAttempts}(a.\mathit{id}) := 0$, $\mathit{sessions} :=
    \mathit{sessions} \cup \{s\}$
  \item When InvalidCredentials is raised,
    $\mathit{failedAttempts}(a.\mathit{id})$ increases by 1. If it then equals
    \textsf{LockoutAttempts}, the account is locked:
    $\mathit{lockedUntil}(a.\mathit{id}) := \mathit{now}() +
    \textsf{LockoutDuration}$ and $\mathit{failedAttempts}(a.\mathit{id}) :=
    0$.
  \end{itemize}
\item output: $\mathit{out} := (\mathit{mfaEnabled}(a.\mathit{id})
  \Rightarrow (\mathit{id}$ of the new challenge$) \mid \neg
  \mathit{mfaEnabled}(a.\mathit{id}) \Rightarrow (\mathit{issueAccess}(s),
  \mathit{rt}, \mathit{now}() + \textsf{AccessExpiry}))$
\item exception, checked in order: $\mathit{overLimit}(\mathit{ip})
  \Rightarrow \mathrm{RateLimited} \mid a = \bot \Rightarrow
  \mathrm{UnknownAccount} \mid \mathit{now}() <
  \mathit{lockedUntil}(a.\mathit{id}) \Rightarrow \mathrm{AccountLocked} \mid
  \neg \mathit{ok} \Rightarrow \mathrm{InvalidCredentials}$
\end{itemize}

\noindent verifyMfa($\mathit{cid}, \mathit{code}, \mathit{ip}$):
\begin{itemize}
\item Let $h$ be the challenge in $\mathit{challenges}$ with $h.\mathit{id} =
  \mathit{cid}$, $\mathit{rt} = \mathit{newSecret}()$, and $s =
  (\mathit{newId}(), h.\mathit{user}, h.\mathit{device},
  \mathit{digest}(\mathit{rt}), \mathit{now}(), \mathrm{false})$.
\item transition: $\mathit{challenges} := \mathit{challenges} - \{h\}$,
  $\mathit{sessions} := \mathit{sessions} \cup \{s\}$
\item output: $\mathit{out} := (\mathit{issueAccess}(s), \mathit{rt},
  \mathit{now}() + \textsf{AccessExpiry})$
\item exception: $\mathit{overLimit}(\mathit{ip}) \Rightarrow
  \mathrm{RateLimited} \mid h$ does not exist $\vee\ \mathit{now}() >
  h.\mathit{expires} \vee \neg \mathit{matches}(\mathit{code},
  h.\mathit{codeHash}) \Rightarrow \mathrm{InvalidCode}$
\end{itemize}

\noindent refresh($\mathit{rt}$):
\begin{itemize}
\item Let $s$ be the session in $\mathit{sessions}$ with $\neg
  s.\mathit{revoked} \wedge s.\mathit{refreshHash} =
  \mathit{digest}(\mathit{rt})$, and $\mathit{rt}' = \mathit{newSecret}()$.
\item transition: $s.\mathit{refreshHash}, s.\mathit{lastUsed} :=
  \mathit{digest}(\mathit{rt}'), \mathit{now}()$. The old refresh token can
  no longer be used.
\item output: $\mathit{out} := (\mathit{issueAccess}(s), \mathit{rt}',
  \mathit{now}() + \textsf{AccessExpiry})$
\item exception: $s$ does not exist $\vee\ \mathit{now}() -
  s.\mathit{lastUsed} > \textsf{SessionTimeout} \Rightarrow
  \mathrm{Unauthenticated}$
\end{itemize}

\noindent authenticate($\mathit{tok}, \mathit{ip}$):
\begin{itemize}
\item Let $p = \mathit{verifyAccess}(\mathit{tok})$ and $s$ be the session in
  $\mathit{sessions}$ with $s.\mathit{id} = p.\mathit{session}$.
\item output: $\mathit{out} := p$
\item exception: $p = \bot \Rightarrow \mathrm{Unauthenticated} \mid
  s.\mathit{revoked} \Rightarrow \mathrm{Unauthenticated} \mid
  \mathit{overLimit}(p.\mathit{user}) \Rightarrow \mathrm{RateLimited}$
\end{itemize}

\noindent authorize($p, \mathit{owner}$):
\begin{itemize}
\item exception: $p.\mathit{user} \neq \mathit{owner} \Rightarrow
  \mathrm{Forbidden}$
\end{itemize}

\noindent requireRole($p, r$):
\begin{itemize}
\item exception: $p.\mathit{role} \neq r \Rightarrow \mathrm{Forbidden}$
\end{itemize}

\noindent confirmPassword($p, \mathit{pw}$):
\begin{itemize}
\item exception: $\neg \mathit{matches}(\mathit{pw},
  \mathit{account}(p.\mathit{user}).\mathit{passwordHash}) \Rightarrow
  \mathrm{InvalidCredentials}$
\end{itemize}

\noindent signOut($p$):
\begin{itemize}
\item transition: $s.\mathit{revoked} := \mathrm{true}$ for the session $s$
  with $s.\mathit{id} = p.\mathit{session}$. Every access and refresh token
  issued for that session, and so for that device, is rejected from then on.
\end{itemize}

\noindent endAllSessions($u$):
\begin{itemize}
\item transition: $s.\mathit{revoked} := \mathrm{true}$ for every $s \in
  \mathit{sessions}$ with $s.\mathit{user} = u$; $\mathit{challenges} :=
  \{h \in \mathit{challenges} \mid h.\mathit{user} \neq u\}$;
  $\mathit{resetTokens} := \{r \in \mathit{resetTokens} \mid r.\mathit{user}
  \neq u\}$
\end{itemize}

\noindent setMfaEnabled($p, \mathit{on}$):
\begin{itemize}
\item transition: $\mathit{mfaEnabled}(p.\mathit{user}) := \mathit{on}$
\end{itemize}

\noindent requestPasswordReset($\mathit{email}, \mathit{ip}$):
\begin{itemize}
\item Let $a = \mathit{findUser}(\mathit{email})$ and $t =
  \mathit{newSecret}()$.
\item transition:
  \begin{itemize}
  \item $a \neq \bot \Rightarrow$ the reset token $(\mathit{digest}(t),
    a.\mathit{id}, \mathit{now}() + \textsf{ResetTokenExpiry})$ is added to
    $\mathit{resetTokens}$, and a link containing $t$ is emailed to
    $a.\mathit{email}$ through $\mathit{mailServer}$
  \item $a = \bot \Rightarrow$ no change
  \end{itemize}
\item output: none in either case
\item exception: $\mathit{overLimit}(\mathit{ip}) \Rightarrow
  \mathrm{RateLimited}$
\end{itemize}

\noindent resetPassword($t, \mathit{pw}$):
\begin{itemize}
\item Let $r$ be the reset token in $\mathit{resetTokens}$ with
  $r.\mathit{tokenHash} = \mathit{digest}(t)$.
\item transition:
  \begin{itemize}
  \item $\mathit{account}(r.\mathit{user}).\mathit{passwordHash} :=
    \mathit{hash}(\mathit{pw})$
  \item $\mathit{failedAttempts}(r.\mathit{user}) := 0$ and
    $\mathit{lockedUntil}(r.\mathit{user}) := 0$, so a locked-out user can
    sign in again
  \item the transition of endAllSessions($r.\mathit{user}$), which also
    removes $r$ so the link cannot be used again
  \end{itemize}
\item exception: $r$ does not exist $\vee\ \mathit{now}() > r.\mathit{expires} \Rightarrow \mathrm{InvalidResetToken}
  \mid \mathit{isWeak}(\mathit{pw}) \Rightarrow \mathrm{WeakPassword}$
\end{itemize}

\subsubsection{Local Functions}

\begin{itemize}
\item UserRecord = tuple of ($\mathit{id}$: UserID, $\mathit{email}$: String,
  $\mathit{role}$: Role, $\mathit{passwordHash}$: String)
\item $\mathit{findUser}$: String $\rightarrow$ UserRecord $\cup \{\bot\}$\\
  $\mathit{findUser}(e) \equiv$ the stored user record whose email is $e$,
  read through the Persistence Module, or $\bot$ if there is none.
\item $\mathit{account}$: UserID $\rightarrow$ UserRecord\\
  $\mathit{account}(u) \equiv$ the stored user record whose id is $u$, read
  and written through the Persistence Module.
\item $\mathit{isWeak}$: String $\rightarrow \mathbb{B}$\\
  $\mathit{isWeak}(\mathit{pw}) \equiv |\mathit{pw}| <
  \textsf{MinPasswordLength} \vee \mathit{pw} \in \mathit{CommonPasswords}$
\item $\mathit{hash}$: String $\rightarrow$ String\\
  A salted one-way hash using a new random salt on each call. It is used for
  passwords and multi-factor authentication codes. (Implemented with BCrypt
  through Spring Security's \texttt{PasswordEncoder}.)
\item $\mathit{matches}$: String $\times$ String $\rightarrow \mathbb{B}$\\
  $\mathit{matches}(x, h) \equiv h$ is the result of $\mathit{hash}(x)$ under
  the salt stored in $h$.
\item $\mathit{digest}$: String $\rightarrow$ String\\
  An unsalted one-way hash (SHA-256). It is used only for long random values
  produced by $\mathit{newSecret}$, so that a stored value can be looked up
  from the token the client presents.
\item $\mathit{newSecret}$: $\rightarrow$ Token\\
  A random value of at least 256 bits from a cryptographically secure random
  number generator.
\item $\mathit{newCode}$: $\rightarrow$ String\\
  A random six-digit code from a cryptographically secure random number
  generator.
\item $\mathit{newId}$: $\rightarrow \mathbb{N}$\\
  An identifier not used by any existing session or challenge.
\item $\mathit{issueAccess}$: Session $\rightarrow$ Token\\
  $\mathit{issueAccess}(s) \equiv$ a token containing $s.\mathit{user}$,
  $\mathit{account}(s.\mathit{user}).\mathit{role}$, $s.\mathit{id}$,
  $s.\mathit{device}$ and the expiry time $\mathit{now}() +
  \textsf{AccessExpiry}$, signed with $\mathit{signingKey}$. (Implemented as a
  JSON Web Token.)
\item $\mathit{verifyAccess}$: Token $\rightarrow$ Principal $\cup \{\bot\}$\\
  $\mathit{verifyAccess}(t) \equiv \bot$ if $t$ is missing, its signature does
  not match $\mathit{signingKey}$, or its expiry time is not after
  $\mathit{now}()$; otherwise the Principal contained in $t$.
\item $\mathit{overLimit}$: String $\rightarrow \mathbb{B}$\\
  $\mathit{overLimit}(k) \equiv |\langle x \in \mathit{requestLog}(k) \mid x >
  \mathit{now}() - 60 \rangle| > \textsf{RateLimit}$
\item $\mathit{log}$: event $\times$ severity $\times$ String $\times$ (UserID
  $\cup \{\bot\}$) $\rightarrow$ none\\
  Writes one structured JSON entry to $\mathit{securityLog}$ containing the
  severity, a UTC timestamp in ISO 8601 format, the client IP address, the
  event type, and the user ID or ``anonymous''. Passwords, codes and tokens
  are never written.
\end{itemize}

\bibliographystyle {plainnat}
\bibliography {../../../refs/References}

\section{Appendix --- Generative AI Use Disclosure}

\wss{ChatGPT (version 5) was used to brainstorm ideas for the template for this disclose section.}

\subsection{AI Tools Used}

\wss{Give the name of the AI tool(s) used and the version.}

\subsection{How and Where Used}

\wss{Explain what the tool(s) were used for.  Examples include: brainstorming,
explaining a technical concept, reviewing the document, generating or modifying
text, generating or modifying diagrams, identifying errors or omissions or
inconsistencies.}

\wss{For each of the uses, identify which parts of the document were affected.}

\wss{You do not need to report every interaction, but you should disclose the 
uses that materially contributed to your work}

\subsection{Verification of AI-Generated Content}

\wss{\begin{itemize}
    \item Describe how the team checked AI-generated or AI-assisted material
    before including it in the document.
    \item In particular, verify factual claims, technical assertions,
    requirements, calculations, references, and other information for which an
    AI system may produce incorrect or fabricated results.
    \item \textbf{The team is responsible for the accuracy, completeness,
consistency, and appropriateness of the submitted document, regardless of
whether material was generated by a human or an AI system.}
\end{itemize}}

\section{Appendix --- Reflection}

\wss{Not required for CAS 741 projects}

The information in this section will be used to evaluate the team members on the
graduate attribute of Problem Analysis and Design.

\input{../../Reflection.text}

\begin{enumerate}
  \item What went well while writing this deliverable? 
  \item What pain points did you experience during this deliverable, and how
    did you resolve them?
  \item Which of your design decisions stemmed from speaking to your client(s)
  or a proxy (e.g. your peers, stakeholders, potential users)? For those that
  were not, why, and where did they come from?
  \item While creating the design doc, what parts of your other documents (e.g.
  requirements, hazard analysis, etc), it any, needed to be changed, and why?
  \item What are the limitations of your solution?  Put another way, given
  unlimited resources, what could you do to make the project better? (LO\_ProbSolutions)
  \item Give a brief overview of other design solutions you considered.  What
  are the benefits and tradeoffs of those other designs compared with the chosen
  design?  From all the potential options, why did you select the documented design?
  (LO\_Explores)
  \item (After you have implemented another team's module, which means this
  isn't filled in until after the original deadline). What did you learn by
  implementing another team's module? Were all the details you needed in the
  documentation, or did you need to make assumptions, or ask the other team
  questions?  If your team also had another team implement one of your modules,
  what was this experience like?  Are there things in your documentation you
  could have changed to make the process go more smoothly for when an
  ``outsider'' completes some of the implementation?

  Your reflection should be about concepts learned in your SE education, such as modular decomposition, cohesion, coupling, information hiding, etc.  You should critique the quality and utility of the specification you received.

  You should create a PR for your implemented module.  A link to the PR should be part of your reflection.

\end{enumerate}

\end{document}